
\chapter{Vacío de la intersección de una RCG con un lenguaje regular finito}
\label{cap:vacio-interseccion-rcg-regular}

El esquema por intersección del capítulo~\ref{cap:rcg-sat} reduce la
satisfacibilidad de una fórmula a decidir si la intersección de dos
lenguajes es vacía~\cite{fernandezarias2007,aguileralopez2016}.

El paso de las gramáticas libres del contexto y de las gramáticas de
concatenación de rango simples a las RCG no conserva de manera
automática la tratabilidad de ese enfoque. Bertsch y Nederhof prueban
que, para RCG, el vacío de la intersección con un lenguaje
regular arbitrario es indecidible. También prueban que, si el lenguaje
regular es finito y se representa mediante un autómata acíclico, el
problema de decidir si la intersección es no vacía es
NP-completo~\cite{bertsch2001rcg}. En el presente capítulo se da un
algoritmo de fuerza bruta para ese último caso.

Este algoritmo no sustituye el resultado de complejidad de
Bertsch y Nederhof. Se precisa qué procedimiento de decisión se obtiene
cuando el autómata es acíclico y se demuestra su correctitud. Su valor
está en exponer qué ocurre al decidir el vacío con los dos formalismos
por separado y en caracterizar su costo. Con él se decide el problema,
pero no se obtiene una representación de la intersección que se pueda
estudiar más allá de las propiedades del autómata. Por eso en el
capítulo~\ref{cap:interseccion}
se construye un objeto que, en un solo formalismo, expresa de manera
simbólica\footnote{Es decir, mediante la gramática misma, y no
recorriendo los dos formalismos por separado.} lo que reconocen ambos.

El capítulo se organiza de la siguiente manera. En la
sección~\ref{sec:lenguajes-regulares-finitos} se presentan los lenguajes
regulares finitos. En la sección~\ref{sec:automatas-aciclicos} se
formalizan los autómatas acíclicos y sus caminos aceptantes. En la
sección~\ref{sec:problema-vacio-interseccion} se define el problema de
vacío de la intersección. Por último, en la
sección~\ref{sec:algoritmo-decision} se presenta el algoritmo de
decisión, se demuestra su correctitud y se analiza su complejidad.

\section{Lenguajes regulares finitos}
\label{sec:lenguajes-regulares-finitos}

El enfoque SAT por intersección trabaja con un lenguaje regular que
representa interpretaciones de una fórmula. Para una fórmula con una
cantidad finita de ocurrencias de variables booleanas, ese lenguaje es finito. La presente
sección fija la noción de lenguaje regular finito que será usada en el
resto del capítulo.

\begin{definicion}
Un lenguaje $L$ es \textbf{finito} si existe un entero $k \ge 0$ tal
que $L$ tiene exactamente $k$ cadenas.
\end{definicion}

\begin{definicion}
Un \textbf{lenguaje regular finito} es un lenguaje regular que también
es finito.
\end{definicion}

Una familia de lenguajes finitos puede tener
una representación mediante autómatas de tamaño polinomial y contener,
al mismo tiempo, una cantidad exponencial de cadenas.

Por ejemplo, para cada $n \ge 1$, sea
\[
  L_n = \{0,1\}^n.
\]
El lenguaje $L_n$ es finito y contiene $2^n$ cadenas. No obstante,
existe un autómata finito con $n+1$ estados que reconoce
$L_n$: el autómata avanza de un estado al siguiente al leer $0$ o $1$, y
acepta solo después de leer exactamente $n$ símbolos.

El ejemplo anterior muestra por qué no basta con saber que el lenguaje
regular es finito. La finitud garantiza que una enumeración termina, pero
no garantiza que termine en tiempo polinomial respecto al tamaño de la
representación. Cuando un lenguaje regular es finito, admite una
representación mediante un autómata acíclico, que se define en la
sección siguiente.

\section{Autómatas acíclicos}
\label{sec:automatas-aciclicos}

En el presente capítulo se usa el autómata finito del
capítulo~\ref{cap:preliminares}, con una diferencia: la función de
transición es parcial, no total, es decir, $\delta(q,a)$ puede no estar
definida para algunos pares $(q,a)$. Cuando $\delta(q,a)$ no está
definida, la lectura del símbolo $a$ en el estado $q$ no conduce a
ningún estado.

\begin{definicion}
Sean $p_0,p_1,\ldots,p_m$ estados de $Q$ y $a_1,a_2,\ldots,a_m$
símbolos de $\Sigma$, con $m \ge 0$. Un \textbf{camino} desde $p_0$
hasta $p_m$ es una secuencia
\[
  \pi = p_0 a_1 p_1 a_2 p_2 \cdots a_m p_m
\]
tal que, para cada $1 \le i \le m$, se cumple
\[
  \delta(p_{i-1},a_i)=p_i.
\]
La cadena que se lee a lo largo de $\pi$, denotada $w(\pi)$, es
$a_1a_2\cdots a_m$. Si $m=0$, el camino no contiene transiciones y la
cadena que se lee es $\varepsilon$.
\end{definicion}

\begin{definicion}
Un camino $\pi$ en el autómata es \textbf{aceptante} si empieza en el
estado inicial $q_0$ y termina en un estado de aceptación. El lenguaje
aceptado por $A$ es
\[
  L(A) =
  \{w(\pi) \mid \pi
  \text{ es un camino aceptante de } A\}.
\]
\end{definicion}

La aciclicidad se define a partir de los caminos del autómata.

\begin{definicion}
Un autómata finito $A$ es \textbf{acíclico} si
ningún camino de $A$ con al menos una transición empieza y termina en
el mismo estado.
\end{definicion}

La aciclicidad es suficiente para asegurar que el lenguaje aceptado es
finito. En los tres lemas siguientes se enuncian la cota de longitud de
los caminos, que se usa en el análisis de complejidad, la finitud del
lenguaje aceptado y la representación de todo lenguaje finito mediante
un autómata acíclico.

\begin{lema}
\label{lema:longitud-caminos-aciclicos}
Sea $A=(Q,\Sigma,\delta,q_0,F)$ un autómata finito acíclico. Todo
camino de $A$ tiene longitud a lo sumo $|Q|-1$. En
consecuencia, toda cadena aceptada por $A$ tiene longitud a lo sumo
$|Q|-1$~\cite{bertsch2001rcg}.
\end{lema}

\begin{lema}
\label{lema:lenguaje-finito-aciclico}
Si $A$ es un autómata finito acíclico, entonces
$L(A)$ es finito~\cite{bertsch2001rcg}.
\end{lema}

\begin{lema}
\label{lema:finito-reconocible-aciclico}
Todo lenguaje finito puede reconocerse mediante un autómata finito
acíclico~\cite{bertsch2001rcg}.
\end{lema}

En el análisis de complejidad del presente capítulo, el autómata forma
parte de la entrada, de modo que el costo se expresa en función de su
tamaño. Ese costo depende, ante todo, de la cantidad de caminos
aceptantes del autómata. Por ello se introduce la siguiente notación.

\begin{definicion}
Sea $A$ un autómata finito acíclico. Se denota por $p(A)$ la cantidad
de caminos aceptantes de $A$.
\end{definicion}

Como $\delta$ asigna a lo sumo un estado a cada par $(q,a)$, cada
cadena aceptada determina un único camino aceptante. En consecuencia,
\[
  p(A)=|L(A)|.
\]

El autómata acíclico es la parte regular de la entrada; en la próxima
sección se define el problema del vacío de la intersección.

\section{Problema de vacío de la intersección}
\label{sec:problema-vacio-interseccion}

El
\textbf{problema del vacío de la intersección de una RCG con un lenguaje
regular finito} toma como entrada una RCG $G$ y un autómata
finito acíclico $A$, y pregunta si $L(G) \cap L(A) \neq \emptyset$.

La intersección es no vacía si y solo si existe una cadena que se
reconoce a la vez por la RCG $G$ y por el autómata $A$:
\[
  L(G) \cap L(A) \neq \emptyset
  \quad \Longleftrightarrow \quad
  \exists w \in \Sigma^* \text{ tal que } w \in L(G) \text{ y } w \in L(A).
\]

Decidir ese problema es el objetivo de la sección siguiente.

\section{Algoritmo de decisión}
\label{sec:algoritmo-decision}

En esta sección se presenta el algoritmo; su correctitud se demuestra en
la subsección~\ref{sec:correctitud-algoritmo} y su complejidad se analiza
en la subsección~\ref{sec:complejidad-algoritmo}.

Como el autómata de entrada es acíclico, sus caminos aceptantes son
finitos. Además, como $\delta$ asigna a lo sumo un estado a cada par
$(q,a)$, esos caminos están en correspondencia uno a uno con las
cadenas aceptadas. Por tanto, una
forma directa de decidir si la intersección es no vacía consiste en
recorrer las cadenas aceptadas por el autómata y aplicar, a cada una, un
algoritmo de reconocimiento para RCG. El procedimiento es
determinista y, en general, no se ejecuta en tiempo polinomial.

Se usará como subrutina un procedimiento de reconocimiento para RCG, que
se denota
\[
  \proc{RCG-Recognize}(G,w).
\]
La subrutina devuelve verdadero si y solo si $w \in L(G)$. El
reconocimiento de una cadena por una RCG se resuelve en tiempo
polinomial en la longitud de la
cadena~\cite{boullier1998proposal,kallmeyer2010parsing}.

Antes de enumerar cadenas conviene eliminar los estados que no pueden
participar en ningún camino aceptante. Esa reducción evita recorrer
ramas que no contribuyen al lenguaje del autómata.

\begin{definicion}
Sea $A=(Q,\Sigma,\delta,q_0,F)$ un autómata finito. Un estado $q\in Q$
es \textbf{útil} si es alcanzable desde $q_0$ y desde
$q$ se puede alcanzar algún estado de aceptación.
\end{definicion}

El conjunto de estados útiles se obtiene mediante dos recorridos del
autómata. En el primero se determinan los estados alcanzables desde
$q_0$: se parte de $\{q_0\}$ y se agrega $\delta(q,a)$ por cada estado
$q$ ya incluido y cada símbolo $a$ para el que $\delta(q,a)$ está
definida, hasta que no se agregue ningún estado. En el segundo se
determinan los estados desde los que se alcanza un estado de aceptación:
se parte de $F$ y se agrega cada estado $q$ tal que $\delta(q,a)$ ya
está incluido para algún símbolo $a$, hasta que no se agregue ninguno.
Un estado es útil cuando pertenece a los dos conjuntos. El autómata
$\proc{Useful-Part}(A)$ tiene por estados los útiles, por estado inicial
$q_0$, por estados de aceptación los de $F$ que son útiles y por
transiciones las $\delta(q,a)=q'$ en las que $q$ y $q'$ son útiles. En
cada recorrido, cada estado y cada transición se visitan a lo sumo una
vez, de modo que la reducción se calcula en tiempo lineal en el tamaño
del autómata.

El pseudocódigo de $\proc{Useful-Part}$ es el siguiente.

\begin{codebox}
\Procname{$\proc{Useful-Part}(A)$}
\li $R = \{q_0\}$
\li \While there is a transition $\delta(q,a)=q'$ with $q \in R$ and $q' \notin R$
\li \Do $R = R \cup \{q'\}$
    \End
\li $C = F$
\li \While there is a transition $\delta(q,a)=q'$ with $q' \in C$ and $q \notin C$
\li \Do $C = C \cup \{q\}$
    \End
\li $U = R \cap C$
\li \Return $(U, \Sigma, \delta', q_0, F \cap U)$, where $\delta'$ is the
    restriction of $\delta$ to $U$
\end{codebox}

El autómata $\proc{Useful-Part}(A)$ acepta el mismo lenguaje que $A$;
cuando no hay estados útiles, ese lenguaje es vacío.

El algoritmo principal se presenta a continuación.

\begin{codebox}
\Procname{$\proc{Nonempty-Intersection}(G,A)$}
\li $A' = \proc{Useful-Part}(A)$
\li \If the initial state of $A'$ does not exist
\li \Then \Return \const{false}
    \End
\li \Return $\proc{Search}(q_0,\varepsilon)$
\end{codebox}

La subrutina $\proc{Search}$ recorre las cadenas aceptadas por el
autómata. El parámetro $q$ indica el estado actual y $w$ almacena la
cadena que se lee a lo largo del camino desde $q_0$ hasta ese estado.

\begin{codebox}
\Procname{$\proc{Search}(q,w)$}
\li \If $q \in F$ \kw{and} $\proc{RCG-Recognize}(G,w)$
\li \Then \Return \const{true}
    \End
\li \For each symbol $a \in \Sigma$ such that $\delta'(q,a)$ is defined
\li \Do
        $r = \delta'(q,a)$
\li     \If $\proc{Search}(r,wa)$
\li     \Then \Return \const{true}
        \End
    \End
\li \Return \const{false}
\end{codebox}

En el pseudocódigo, $\delta'$ denota la función de transición de $A'$.
La línea~1 de $\proc{Search}$ verifica la pertenencia a la RCG cada vez
que el camino recorrido termina en un estado de aceptación. Si la
verificación
es positiva, la intersección es no vacía y el algoritmo termina. Si la
verificación es negativa, el recorrido continúa, ya que desde un estado
de aceptación pueden salir transiciones hacia otros caminos aceptantes.
La línea
5 actualiza la cadena asociada al camino mediante la concatenación del
símbolo leído.

La aciclicidad de $A$ garantiza la terminación de la búsqueda. Los
estados $q_0,q_1,\ldots,q_d$ que aparecen en una secuencia de llamadas
recursivas anidadas forman un camino de $A'$: en cada paso, de la
llamada en el estado $q$ se pasa a la llamada en el estado
$\delta'(q,a)$. La profundidad $d$ de la recursión es entonces el número
de transiciones de ese camino. Como $A'$ es acíclico, ese camino no
contiene dos veces el mismo estado, de modo que ningún estado se repite
a lo largo de la recursión. Por el
lema~\ref{lema:longitud-caminos-aciclicos}, $d \le |Q|-1$. Como en cada
llamada se generan a lo sumo $|\Sigma|$ llamadas recursivas, el árbol de
recursión es finito y la búsqueda termina.

A continuación se demuestra la correctitud del algoritmo.

\subsection{Correctitud}
\label{sec:correctitud-algoritmo}

La demostración separa tres hechos. El primero establece
que toda cadena que se evalúa con $\proc{RCG-Recognize}$ pertenece al
lenguaje del autómata. El segundo establece que toda cadena aceptada por
el autómata se evalúa exactamente una vez. El tercero se apoya en la
correctitud de $\proc{RCG-Recognize}$.

\begin{lema}
\label{lema:sonido-automata}
Toda cadena $w$ para la que el algoritmo ejecuta
$\proc{RCG-Recognize}(G,w)$ pertenece a $L(A)$.
\end{lema}

\begin{proof}
La subrutina $\proc{Search}$ ejecuta $\proc{RCG-Recognize}(G,w)$ solo en
la línea~1, cuando el estado actual $q$ pertenece a $F$. El valor de $w$
se construye concatenando los símbolos que se leen desde $q_0$ hasta
$q$. Por tanto, existe un camino desde $q_0$ hasta un estado de
aceptación a lo
largo del cual se lee $w$. Por la definición de aceptación,
$w \in L(A')$. Como $A'=\proc{Useful-Part}(A)$ acepta el mismo lenguaje
que $A$, se tiene $L(A')=L(A)$. Luego $w \in L(A)$.
\end{proof}

\begin{lema}
\label{lema:completo-automata}
Para toda cadena $w \in L(A)$, el algoritmo ejecuta
$\proc{RCG-Recognize}(G,w)$ exactamente una vez.
\end{lema}

\begin{proof}
Sea $w=a_1a_2\cdots a_m$ una cadena de $L(A)$. Como para cada par
$(q,a)$ hay a lo sumo un estado $\delta(q,a)$, existe un único camino de
$A$ que comienza en $q_0$ y a lo largo del cual se lee $w$:
\[
  q_0 a_1 q_1 a_2 q_2 \cdots a_m q_m.
\]
Además, como $w \in L(A)$, se tiene $q_m \in F$. Todos los estados de ese
camino son útiles: cada uno es alcanzable desde $q_0$ por el prefijo
correspondiente, y desde cada uno se alcanza un estado de aceptación
por el
sufijo correspondiente. Por tanto, el camino también pertenece a $A'$.

La subrutina $\proc{Search}$ recorre, para cada estado actual, todos los
símbolos $a$ para los cuales $\delta'(q,a)$ está definida. Por inducción
sobre la longitud de los prefijos de $w$, el algoritmo realiza una
llamada recursiva asociada al único prefijo de longitud $i$ de $w$, para
cada $0 \le i \le m$. Al llegar al estado $q_m$, la cadena acumulada es
$w$ y se cumple $q_m \in F$. En ese momento se ejecuta
$\proc{RCG-Recognize}(G,w)$.

La ejecución es única porque en $A'$ hay a lo sumo un estado
$\delta'(q,a)$ por cada par $(q,a)$: ninguna cadena se lee a lo largo
de dos caminos distintos desde $q_0$.
Por tanto, el algoritmo no puede alcanzar un estado de aceptación con
la misma
cadena $w$ por dos recorridos diferentes.
\end{proof}

\begin{lema}
\label{lema:correctitud-reconocedor}
Para toda cadena $w$, la subrutina $\proc{RCG-Recognize}(G,w)$ devuelve
verdadero si y solo si $w \in L(G)$.
\end{lema}

\begin{proof}
El reconocimiento de una cadena por una RCG se decide en tiempo
polinomial~\cite{boullier1998proposal,kallmeyer2010parsing}. $\proc{RCG-Recognize}$
es un procedimiento que resuelve ese problema y devuelve verdadero si
y solo si $w \in L(G)$.
\end{proof}

\begin{teorema}
\label{teo:correctitud-enumeracion}
El algoritmo $\proc{Nonempty-Intersection}(G,A)$ devuelve
verdadero si y solo si
\[
  L(G) \cap L(A) \neq \emptyset.
\]
\end{teorema}

\begin{proof}
Supóngase primero que el algoritmo devuelve verdadero. Entonces, para
alguna cadena $w$, se ejecutó $\proc{RCG-Recognize}(G,w)$ y la subrutina
devolvió verdadero. Por el lema~\ref{lema:sonido-automata},
$w \in L(A)$. Por el lema~\ref{lema:correctitud-reconocedor},
$w \in L(G)$. Luego $w \in L(G) \cap L(A)$, y la intersección es no
vacía.

Supóngase ahora que $L(G) \cap L(A) \neq \emptyset$. Entonces existe una
cadena $w$ tal que $w \in L(G)$ y $w \in L(A)$. Por el
lema~\ref{lema:completo-automata}, el algoritmo ejecuta
$\proc{RCG-Recognize}(G,w)$ exactamente una vez. Por el
lema~\ref{lema:correctitud-reconocedor}, esa ejecución devuelve
verdadero. Por tanto, el algoritmo devuelve verdadero.
\end{proof}

A continuación se analiza la complejidad del algoritmo.

\subsection{Complejidad}
\label{sec:complejidad-algoritmo}

El análisis depende de dos componentes. El primero es
la enumeración de cadenas aceptadas por el autómata acíclico. El
segundo es el reconocimiento de cada cadena mediante
la RCG. En la presente subsección se expresan ambos costos de manera
separada.

Sea $|A|=|Q|+|\delta|$ el tamaño del autómata, donde $|\delta|$ denota
la cantidad de pares $(q,a)$ para los cuales la transición está definida.
Sea $p(A)$ el número de caminos aceptantes de $A$. Como se estableció
en la sección~\ref{sec:automatas-aciclicos}, $p(A)=|L(A)|$.

Después de eliminar estados no útiles, todo camino recorrido por el
algoritmo es prefijo de algún camino aceptante. Como cada camino
aceptante tiene longitud a lo sumo $|Q|-1$, el número de prefijos
recorridos está acotado por
\[
  p(A)\,|Q|.
\]
Por tanto, el costo de la enumeración, sin contar las llamadas a
$\proc{RCG-Recognize}$, queda acotado por
\[
  O(p(A)\,|Q|).
\]

Sea $R_G(n)$ el tiempo que demora
$\proc{RCG-Recognize}(G,w)$ sobre una cadena $w$ de longitud $n$. Para
una RCG fija, $R_G(n)$ es polinomial en $n$. Si la gramática forma parte
de la entrada, el costo es polinomial en la longitud de la cadena para
cada gramática fijada, y el exponente depende de parámetros
estructurales de la gramática, como la cantidad de variables que pueden
recibir rangos en una cláusula~\cite{boullier1998proposal,kallmeyer2010parsing}.

Por el lema~\ref{lema:longitud-caminos-aciclicos}, toda cadena aceptada
tiene longitud a lo sumo $|Q|-1$. Luego el costo total del algoritmo se
acota por
\[
  O\bigl( |A| + p(A)\,|Q| + p(A)\,R_G(|Q|-1) \bigr).
\]
De forma equivalente,
\[
  O\bigl( |A| + p(A)(|Q| + R_G(|Q|-1)) \bigr).
\]

La cota anterior es polinomial si $p(A)$ está acotado por un polinomio
en $|A|$. Sin embargo, $p(A)$ no está polinomialmente acotado para todos
los autómatas acíclicos. En la
sección~\ref{sec:lenguajes-regulares-finitos} ya se mostró un
autómata así: el que reconoce $L_n=\{0,1\}^n$, con $n+1$ estados,
acíclico y con $2^n$ cadenas aceptadas. Por la igualdad $p(A)=|L(A)|$,
ese autómata tiene $2^n$ caminos aceptantes con un tamaño lineal en $n$.

Por tanto, la enumeración determinista decide el problema, pero su
costo puede ser exponencial en el tamaño del autómata, incluso para
autómatas acíclicos.

El costo del algoritmo depende de cuántas cadenas acepta la
representación regular. Si una clase de autómatas acíclicos tiene
polinomialmente muchas cadenas aceptadas, el vacío de la intersección
se decide en tiempo polinomial para cada RCG fija; si admite autómatas
con exponencialmente muchas cadenas, ese costo es exponencial. Esa
condición es trivial: enumerar es barato cuando hay poco que enumerar,
de modo que la enumeración no es una vía factible para decidir el
vacío de la intersección, pero sí necesaria para entender el problema
y por qué conviene resolverlo sobre un único objeto que exprese lo que
reconocen ambos formalismos.

Generar cada cadena aceptada en tiempo polinomial no es lo mismo que
enumerarlas todas en tiempo polinomial. La cantidad de cadenas
aceptadas puede ser exponencial, de modo que recorrerlas todas cuesta
tiempo exponencial aun cuando cada una se obtenga en tiempo polinomial.
Para decidir el vacío de la intersección hace falta, en el peor caso,
recorrer todas: cuando la intersección es vacía, ninguna cadena aceptada
pertenece a $L(G)$, y el recorrido solo concluye al agotarlas.

En este capítulo se definió el problema del vacío de la intersección de
una RCG con un lenguaje regular finito y se dio un algoritmo de fuerza
bruta que lo decide cuando el lenguaje regular se representa con un
autómata acíclico; se demostró su correctitud y se caracterizó su costo,
exponencial en el peor caso. En el próximo capítulo se modela el SAT
como ese vacío, asociando a cada fórmula booleana un autómata booleano y
una RCG de consistencia.

